\documentclass{EQDclass}
%\documentclass[blind]{EQDclass}% Use the blind option to produce a blind copy of the manuscript

% Add here your packages
% \usepackage{mypackage}

% Add here your LaTeX macros
% \def\mymacro{}

% Load your bib source file here
\addbibresource{main.bib}
\begin{document}
% Change the title with the title of your paper
\Title{My title}
% Write the list of authors in order
\Author[Title, Affiliation \autemail{myemail@mydomain.com}]{Jane Doe}
%\Author[Title, Affiliation \autemail{hisemail@mydomain.com}]{John Doe}% you can add more authors
% Write your abstract
\Abstract{This is the abstract}
% Write the keywords (they won't appear in the pdf)
\Keywords{Sample, LaTeX}
% Write your JEL classifications code (they won't appear in the pdf)
\JEL{Xnn, Xnn}

% Print the title
\Maketitle


% What follows is an example of elements you can use
\section{Review our submission guidelines}

Prior to submitting your paper, please review our submission guidelines here: \url{https://eqderivatives.com/journal-of-systematic-investing/jsi-submission-guidelines}

\section{This is an example of a section}

This is a regular paragraph. Please do not introduce personal settings to the layout of the section.

\subsection{This is an example of a subsection}

This is a regular paragraph. Don't use any section lower than subsection. If you need a footnote.\footnote{Here's your footnote.}

\section{Some floats}

\subsection{Some tables}

Check the code to introduce a table.

\Cref{Tab:1} is a one-column table, \Cref{Tab:2} is a two-column table.

\begin{table}[!ht]\centering
    \caption{One-column table}\label{Tab:1}
    \begin{tabular}{|c|c|c|}
        \hline
        \TABhead Head 1             & \TABhead Head 2       & \TABhead Head 3       \\ \hline
        \TABsubhead Subhead 1       & \TABsubhead Subhead 2 & \TABsubhead Subhead 3 \\ \hline
        \TABsubsubhead Subsubhead 1 & a                     & b                     \\ \hline
        \TABsubsubhead Subsubhead 2 & a                     & b                     \\ \hline
        \TABsubsubhead Subsubhead 3 & a                     & b                     \\ \hline
        \TABsubsubhead Subsubhead 4 & a                     & b                     \\ \hline
    \end{tabular}
    \afternote{Use this macro to write notes outside the caption.}
    \afternotesource{Use this macro if you need to write the source.}
\end{table}

And here's some random text.

\begin{table*}[!ht]
    \caption{Two-column table}\label{Tab:2}
    \begin{tabular}{|c|c|c|c|c|c|}
        \hline
        \TABhead Head 1             & \TABhead Head 2       & \TABhead Head 3       & \TABhead Head 4       & \TABhead Head 5       & \TABhead Head 6       \\ \hline
        \TABsubhead Subhead 1       & \TABsubhead Subhead 2 & \TABsubhead Subhead 3 & \TABsubhead Subhead 4 & \TABsubhead Subhead 5 & \TABsubhead Subhead 6 \\ \hline
        \TABsubsubhead Subsubhead 1 & \TABsubsubhead 2      & \TABsubsubhead 3      & \TABsubsubhead 4      & \TABsubsubhead 5      & \TABsubsubhead 6      \\ \hline
        \TABsubsubhead Subsubhead 2 & a                     & b                     & c                     & d                     & e                     \\ \hline
        \TABsubsubhead Subsubhead 3 & a                     & b                     & c                     & d                     & e                     \\ \hline
        \TABsubsubhead Subsubhead 4 & a                     & b                     & c                     & d                     & e                     \\ \hline
    \end{tabular}
\end{table*}

And here's some more random text.


\subsection{Some figures}

Check the code to introduce a figure.

Here are some figure examples. \Crefs{Fig:1} and \ref{Fig:2} are two figures. \Cref{Fig:1} is a one-column figure, \Cref{Fig:2} is a two-column figure.

\begin{figure}[!ht]\centering
    \includegraphics[width=\linewidth]{figure1}
    \caption{One-column figure}\label{Fig:1}
\end{figure}

And here's some random text.

\begin{figure*}[!ht]\centering
    \includegraphics[width=0.8\linewidth]{figure2}
    \caption{Two-column figure}\label{Fig:2}
\end{figure*}

And here's some more random text.


\section{Some math}

You can write math using the regular macros in the \verb+amsmath+ package. Here some math examples:
\begin{align}
    \nabla \cdot \mathbf{E}  & = \dfrac{\rho}{\varepsilon_{0}} \label{Eq:1}\\
    \nabla \cdot \mathbf{B}  & = 0 \\
    \nabla \times \mathbf{E} & = -\dfrac{\partial \mathbf{B}}{\partial t} \\
    \nabla \times \mathbf{B} & = \mu_{0} \left(\mathbf{J}+\varepsilon_{0}\dfrac{\partial\mathbf{E}}{\partial t}\right) \label{Eq:4}
\end{align}
Some inline $\mathrm{E}=m\mathrm{c}^2$ math. See formula~\EQref{Eq:1} and \EQrefs{Eq:4}.


\section{Lists}

You can use the classic \LaTeX\ lists environments.


\subsection{Itemize}

\begin{itemize}
    \item Item 1
    \item Item 2
        \begin{itemize}
            \item item a
            \item item b
            \item item c
        \end{itemize}
    \item Item 3
\end{itemize}


\subsection{Enumerate}

\begin{enumerate}
    \item Item 1
    \item Item 2
        \begin{enumerate}
            \item item a
            \item item b
            \item item c
        \end{enumerate}
    \item Item 3
\end{enumerate}


\subsection{Description}

\begin{description}
    \item[Description 1] Item 1
    \item[Description 2] Item 2
    \item[Description 3] Item 3
\end{description}

\section{References}

The JSI adhers to the Chicago Manual of Style, please refer to their guide for more: \url{https://www.chicagomanualofstyle.org/home.html}

To write an in-line citation you can use \verb+\citet+ like this: \citet{latexcompanion} say \LaTeX\ is great.

To build a citation within parenthesis you can use \verb+\citep+ \citep{latexcompanion}.

For multiple citations, you can include them all inside the cite macro \citep{latexcompanion,einstein}.


% Print your footnotes at the end here
\printendnotes

% Appendix
\appendix
\section{Appendix}

This the appendix section, if you need one.


% Load your bibliography
\nocite{knuthwebsite}% use nocite if you want to include the entries from your bib file with no citation in the text. Do not use \nocite{*}
\printbibliography

% Acknowledgments
\begin{blind}
\section{Acknowledgments}

You can write your Acknowledgments here.
If you use the blind environment to write material, it won't be printed with the blind option on.
\end{blind}

% Disclosure
\begin{blind}
\section{Disclosure}

You can write your Disclosure here.
\end{blind}

\end{document}